% Catalogue of \paragraph heading designs for the internship report.
% Compile from the project root: latexmk -pdf paragraph-styles-showcase.tex

\documentclass[en]{telereport}           % same class as the report: same fonts, header, footer
% ---------------------------------------------------------------- extra packages
% The class already loads fontenc, babel, geometry, graphicx, tikz, hyperref...
% Anything else the report needs goes here.

\usepackage{booktabs}            % \toprule, \midrule, \bottomrule for clean tables
\usepackage{amsmath}             % equation environment, \qquad, math spacing
\usepackage{amssymb}             % \mathbb{R} and other symbols (affine registration)
\usepackage{xurl}                % let long URLs break across lines
\usepackage[round]{natbib}       % \cite{key} -> (Zhang et al., 2023), author-year style
\bibliographystyle{apalike-nd}   % APA-like layout (local copy of apalike: keeps "n.d." in citations)
\let\cite\citep                 % \cite{key} -> (Zhang et al., 2023); use \citet for "Zhang et al. (2023)"
\renewcommand{\bibsection}{%           % reference list heading:
  \section*{List of References}%        %   never numbered, even when sections are
  \addcontentsline{toc}{section}{List of References}%  %   but still listed in the table of contents
  \markboth{List of References}{}}      %   and shown in the page header

% ---- code listings (used in the Registration section)
\usepackage{listings}            % typeset source code blocks
\lstset{%
  language=Python,               % keyword highlighting for Python
  basicstyle=\ttfamily\small,    % monospaced, slightly smaller than the text
  keywordstyle=\bfseries,        % keywords in bold
  stringstyle=\color{black!70},  % strings in dark grey
  frame=single,                  % thin box around the code
  framesep=4pt,                  % space between box and code
  xleftmargin=6pt, xrightmargin=6pt,  % keep the box inside the text width
  columns=fullflexible,          % natural spacing, copy-paste friendly
  aboveskip=0.8em, belowskip=0.8em}   % air around each block

% ---- TikZ libraries (used for the registration flowchart)
\usetikzlibrary{positioning,fit,arrows.meta}  % node placement, grouping boxes, arrow tips

% ---- monospaced font: Latin Modern Mono (vector) when installed, e.g. on Overleaf
\IfFileExists{t1lmtt.fd}{\renewcommand{\ttdefault}{lmtt}}{}  % otherwise keep the default

\usepackage{placeins}                  % \FloatBarrier: print pending figures before going on                   % booktabs, amsmath, listings...
% ------------------------------------------------------------------ paragraph heading styles
% Every \paragraph{...} heading in the report is drawn by ONE style, chosen in main.tex with
%     \setparagraphstyle{<name>}
% Styles: tricolour, tricolour-title, accent, fade, tag, spaced, flag, minimal, original.
% Write titles WITHOUT a final full stop: the run-in styles add their own separator.
% Colours come from the class: tpgrey, black and tpred (the footer bar).

\usepackage{microtype}                   % \textls{} = letter spacing, used by the "spaced" style
\makeatletter

% ---- user setting
\newlength{\paragraphtextgap}          % space between the tricolour underline and the text below it
\setlength{\paragraphtextgap}{0.9em}   % default; change it in main.tex with \setlength

% ---- building blocks
\newcommand{\tp@lineunder}[1]{%        % put #1 (a line) right under the title, no extra blank line
  \par\nointerlineskip\vspace{0.45ex}% % small fixed gap instead of a full line of spacing
  \hbox{#1}}                            % the line itself, as a box in the vertical list
\newcommand{\tp@tricolourrule}[2]{%      % #1 total width, #2 thickness: grey | black | red, equal thirds
  {\color{tpgrey}\rule{\dimexpr#1/3\relax}{#2}}%  % first third
  {\color{black}\rule{\dimexpr#1/3\relax}{#2}}%   % second third
  {\color{tpred}\rule{\dimexpr#1/3\relax}{#2}}}   % last third
\newcommand{\tp@miniflag}{%              % tiny footer-like flag for run-in headings
  \raisebox{0.32ex}{\tp@tricolourrule{16pt}{2.6pt}}}  % sits at mid x-height
\newcommand{\tp@titleunderrule}[1]{%     % title with a tricolour rule exactly as wide as the title
  \sbox\@tempboxa{#1}%                   % measure the title
  \vtop{\offinterlineskip\copy\@tempboxa\kern0.45ex%  % title on the baseline, small fixed gap below
        \hbox{\tp@tricolourrule{\wd\@tempboxa}{1.2pt}}}}  % rule of the same width
\newcommand{\tp@tagtitle}[1]{%           % title inside a light label with a red left edge
  \tikz[baseline=(t.base)]{%             % align the label with the text baseline
    \node[fill=tpgrey!10, inner xsep=7pt, inner ysep=3.5pt, font=\small\bfseries] (t) {#1};  % label
    \fill[tpred] (t.north west) rectangle ([xshift=2.2pt]t.south west);}}  % red edge on the left
\newcommand{\tp@spacedtitle}[1]{%        % letter-spaced red capitals followed by a grey line
  {\color{tpred}\small\bfseries\textls[110]{\MakeUppercase{#1}}}%  % the title
  \hspace{0.8em}{\color{tpgrey!45}\leaders\hrule height 2.9pt depth -2.5pt\hfill}}  % line to the margin

% ---- the styles (each one redefines \paragraph with titlesec)
\newcommand{\tp@defstyle}[2]{%        % register a style under any name, hyphens allowed
  \expandafter\def\csname tp@pstyle@#1\endcsname{#2}}
\tp@defstyle{tricolour}{%      % YOUR IDEA: title above a full-width tricolour line
  \titleformat{\paragraph}[block]{\normalfont\normalsize\bfseries}{}{0pt}{}%
    [\tp@lineunder{\tp@tricolourrule{\linewidth}{0.8pt}}]%
  \titlespacing*{\paragraph}{0pt}{1.4em}{0.25em}}
\tp@defstyle{tricolour-title}{% title underlined by a tricolour rule as wide as the title
  \titleformat{\paragraph}[block]{\normalfont\normalsize\bfseries}{}{0pt}{\tp@titleunderrule}%
  \titlespacing*{\paragraph}{0pt}{1.4em}{\paragraphtextgap}}  % gap below = the user setting
\tp@defstyle{accent}{%         % title above a grey hairline that starts with a tricolour flag
  \titleformat{\paragraph}[block]{\normalfont\normalsize\bfseries}{}{0pt}{}%
    [\tp@lineunder{\rlap{\tp@tricolourrule{1.5cm}{1.6pt}}%
     {\color{tpgrey!35}\rule[0.6pt]{\linewidth}{0.4pt}}}]%
  \titlespacing*{\paragraph}{0pt}{1.4em}{0.25em}}
\tp@defstyle{fade}{%           % title above a red line that fades out to the right
  \titleformat{\paragraph}[block]{\normalfont\normalsize\bfseries}{}{0pt}{}%
    [\tp@lineunder{\tikz{\shade[left color=tpred, right color=white]
       (0,0) rectangle (\linewidth,1pt);}}]%
  \titlespacing*{\paragraph}{0pt}{1.4em}{0.25em}}
\tp@defstyle{tag}{%            % title in a small label, text below
  \titleformat{\paragraph}[block]{\normalfont}{}{0pt}{\tp@tagtitle}%
  \titlespacing*{\paragraph}{0pt}{1.3em}{0.6em}}
\tp@defstyle{spaced}{%         % red spaced capitals + thin line to the right margin
  \titleformat{\paragraph}[block]{\normalfont}{}{0pt}{\tp@spacedtitle}%
  \titlespacing*{\paragraph}{0pt}{1.3em}{0.45em}}
\tp@defstyle{flag}{%           % run-in: mini tricolour flag + bold title, text continues
  \titleformat{\paragraph}[runin]{\normalfont\normalsize\bfseries}{}{0pt}{\tp@miniflag\hspace{0.45em}}%
  \titlespacing*{\paragraph}{0pt}{1.1em}{0.7em}}
\tp@defstyle{minimal}{%        % run-in: red bold title, text continues
  \titleformat{\paragraph}[runin]{\normalfont\normalsize\bfseries\color{tpred}}{}{0pt}{}%
  \titlespacing*{\paragraph}{0pt}{1.1em}{0.7em}}
\tp@defstyle{original}{%      % LaTeX default look: bold run-in title with a full stop
  \titleformat{\paragraph}[runin]{\normalfont\normalsize\bfseries}{}{0pt}{}[.]%
  \titlespacing*{\paragraph}{0pt}{3.25ex plus 1ex minus .2ex}{1em}}

% ---- the switch
\newcommand{\setparagraphstyle}[1]{%     % \setparagraphstyle{tricolour}, {tag}, ...
  \@ifundefined{tp@pstyle@#1}%           % unknown name -> clear error message
    {\PackageError{paragraph-styles}{Unknown paragraph style '#1'}%
      {Use tricolour, tricolour-title, accent, fade, tag, spaced, flag, minimal or original.}}%
    {\csname tp@pstyle@#1\endcsname}}    % apply the chosen style
\makeatother

\setparagraphstyle{tricolour-title}      % default; main.tex can override it
           % the styles being presented
\usepackage{tcolorbox}                   % frames around the sample texts
\usepackage{courier}                    % vector monospaced font for style names and code

\title{Paragraph heading styles}         % shown in the header
\shorttitle{Paragraph heading styles -- design catalogue}  % header text

\newtcolorbox{sample}{%                  % a "sample page" card for each style
  colback=white, colframe=tpgrey!25, boxrule=0.4pt, arc=0pt,
  left=12pt, right=12pt, top=0pt, bottom=8pt,
  before upper={\csname @minipagefalse\endcsname}}  % let headings inside the frame keep their space above
\newcommand{\sampletext}{%               % real text from the report, identical in every card
  \paragraph{Encoder design}
  The encoder is a stack of 3D convolutional blocks. Each block consists of a
  $3\times3\times3$ convolution without bias, instance normalization with learnable affine
  parameters, and an ELU activation. Downsampling is performed with strided convolutions rather
  than pooling, in every second layer.
  \paragraph{Latent sampling}
  During training, a latent vector is sampled with the reparameterization trick,
  \[\mathbf{z} = \boldsymbol{\mu} + \exp(\log\boldsymbol{\sigma}) \odot \boldsymbol{\epsilon},
    \qquad \boldsymbol{\epsilon} \sim \mathcal{N}(\mathbf{0}, \mathbf{I}),\]
  and at inference time the mean is used directly, giving a deterministic encoding of each
  airway tree.}
\newcommand{\styleinfo}[3]{%             % one line of facts: name to type, layout, report length
  \noindent{\small\color{tpgrey}Name:~\texttt{\textbackslash setparagraphstyle\{#1\}}%
  \quad$\cdot$\quad Layout:~#2\quad$\cdot$\quad Report:~#3~pages}\par\smallskip}
\newcommand{\showstyle}[5]{%             % #1 number, #2 title, #3 key, #4 layout, #5 pages ; description follows
  \subsection*{#1\quad #2}\styleinfo{#3}{#4}{#5}}

\begin{document}

\section*{Paragraph heading styles}      % page 1: what this is and how to use it
\markboth{Catalogue}{}

Your report contains 101 \texttt{\textbackslash paragraph\{\ldots\}} headings. They are all
drawn by a single style, chosen with one line in \texttt{main.tex}; the section files do not
change. The eight designs below are built from the colours the report already uses -- grey,
black and the red of the footer bar -- so they match the cover, the header and the section
markers.

\subsection*{How to switch}
\begin{lstlisting}[language={[LaTeX]TeX}]
% ------------------------------------------------------------------ paragraph heading styles
% Every \paragraph{...} heading in the report is drawn by ONE style, chosen in main.tex with
%     \setparagraphstyle{<name>}
% Styles: tricolour, tricolour-title, accent, fade, tag, spaced, flag, minimal, original.
% Write titles WITHOUT a final full stop: the run-in styles add their own separator.
% Colours come from the class: tpgrey, black and tpred (the footer bar).

\usepackage{microtype}                   % \textls{} = letter spacing, used by the "spaced" style
\makeatletter

% ---- user setting
\newlength{\paragraphtextgap}          % space between the tricolour underline and the text below it
\setlength{\paragraphtextgap}{0.9em}   % default; change it in main.tex with \setlength

% ---- building blocks
\newcommand{\tp@lineunder}[1]{%        % put #1 (a line) right under the title, no extra blank line
  \par\nointerlineskip\vspace{0.45ex}% % small fixed gap instead of a full line of spacing
  \hbox{#1}}                            % the line itself, as a box in the vertical list
\newcommand{\tp@tricolourrule}[2]{%      % #1 total width, #2 thickness: grey | black | red, equal thirds
  {\color{tpgrey}\rule{\dimexpr#1/3\relax}{#2}}%  % first third
  {\color{black}\rule{\dimexpr#1/3\relax}{#2}}%   % second third
  {\color{tpred}\rule{\dimexpr#1/3\relax}{#2}}}   % last third
\newcommand{\tp@miniflag}{%              % tiny footer-like flag for run-in headings
  \raisebox{0.32ex}{\tp@tricolourrule{16pt}{2.6pt}}}  % sits at mid x-height
\newcommand{\tp@titleunderrule}[1]{%     % title with a tricolour rule exactly as wide as the title
  \sbox\@tempboxa{#1}%                   % measure the title
  \vtop{\offinterlineskip\copy\@tempboxa\kern0.45ex%  % title on the baseline, small fixed gap below
        \hbox{\tp@tricolourrule{\wd\@tempboxa}{1.2pt}}}}  % rule of the same width
\newcommand{\tp@tagtitle}[1]{%           % title inside a light label with a red left edge
  \tikz[baseline=(t.base)]{%             % align the label with the text baseline
    \node[fill=tpgrey!10, inner xsep=7pt, inner ysep=3.5pt, font=\small\bfseries] (t) {#1};  % label
    \fill[tpred] (t.north west) rectangle ([xshift=2.2pt]t.south west);}}  % red edge on the left
\newcommand{\tp@spacedtitle}[1]{%        % letter-spaced red capitals followed by a grey line
  {\color{tpred}\small\bfseries\textls[110]{\MakeUppercase{#1}}}%  % the title
  \hspace{0.8em}{\color{tpgrey!45}\leaders\hrule height 2.9pt depth -2.5pt\hfill}}  % line to the margin

% ---- the styles (each one redefines \paragraph with titlesec)
\newcommand{\tp@defstyle}[2]{%        % register a style under any name, hyphens allowed
  \expandafter\def\csname tp@pstyle@#1\endcsname{#2}}
\tp@defstyle{tricolour}{%      % YOUR IDEA: title above a full-width tricolour line
  \titleformat{\paragraph}[block]{\normalfont\normalsize\bfseries}{}{0pt}{}%
    [\tp@lineunder{\tp@tricolourrule{\linewidth}{0.8pt}}]%
  \titlespacing*{\paragraph}{0pt}{1.4em}{0.25em}}
\tp@defstyle{tricolour-title}{% title underlined by a tricolour rule as wide as the title
  \titleformat{\paragraph}[block]{\normalfont\normalsize\bfseries}{}{0pt}{\tp@titleunderrule}%
  \titlespacing*{\paragraph}{0pt}{1.4em}{\paragraphtextgap}}  % gap below = the user setting
\tp@defstyle{accent}{%         % title above a grey hairline that starts with a tricolour flag
  \titleformat{\paragraph}[block]{\normalfont\normalsize\bfseries}{}{0pt}{}%
    [\tp@lineunder{\rlap{\tp@tricolourrule{1.5cm}{1.6pt}}%
     {\color{tpgrey!35}\rule[0.6pt]{\linewidth}{0.4pt}}}]%
  \titlespacing*{\paragraph}{0pt}{1.4em}{0.25em}}
\tp@defstyle{fade}{%           % title above a red line that fades out to the right
  \titleformat{\paragraph}[block]{\normalfont\normalsize\bfseries}{}{0pt}{}%
    [\tp@lineunder{\tikz{\shade[left color=tpred, right color=white]
       (0,0) rectangle (\linewidth,1pt);}}]%
  \titlespacing*{\paragraph}{0pt}{1.4em}{0.25em}}
\tp@defstyle{tag}{%            % title in a small label, text below
  \titleformat{\paragraph}[block]{\normalfont}{}{0pt}{\tp@tagtitle}%
  \titlespacing*{\paragraph}{0pt}{1.3em}{0.6em}}
\tp@defstyle{spaced}{%         % red spaced capitals + thin line to the right margin
  \titleformat{\paragraph}[block]{\normalfont}{}{0pt}{\tp@spacedtitle}%
  \titlespacing*{\paragraph}{0pt}{1.3em}{0.45em}}
\tp@defstyle{flag}{%           % run-in: mini tricolour flag + bold title, text continues
  \titleformat{\paragraph}[runin]{\normalfont\normalsize\bfseries}{}{0pt}{\tp@miniflag\hspace{0.45em}}%
  \titlespacing*{\paragraph}{0pt}{1.1em}{0.7em}}
\tp@defstyle{minimal}{%        % run-in: red bold title, text continues
  \titleformat{\paragraph}[runin]{\normalfont\normalsize\bfseries\color{tpred}}{}{0pt}{}%
  \titlespacing*{\paragraph}{0pt}{1.1em}{0.7em}}
\tp@defstyle{original}{%      % LaTeX default look: bold run-in title with a full stop
  \titleformat{\paragraph}[runin]{\normalfont\normalsize\bfseries}{}{0pt}{}[.]%
  \titlespacing*{\paragraph}{0pt}{3.25ex plus 1ex minus .2ex}{1em}}

% ---- the switch
\newcommand{\setparagraphstyle}[1]{%     % \setparagraphstyle{tricolour}, {tag}, ...
  \@ifundefined{tp@pstyle@#1}%           % unknown name -> clear error message
    {\PackageError{paragraph-styles}{Unknown paragraph style '#1'}%
      {Use tricolour, tricolour-title, accent, fade, tag, spaced, flag, minimal or original.}}%
    {\csname tp@pstyle@#1\endcsname}}    % apply the chosen style
\makeatother

\setparagraphstyle{tricolour-title}      % default; main.tex can override it
   % loads all the designs
\setparagraphstyle{tricolour}    % pick one: tricolour, tricolour-title,
                                 % accent, fade, tag, spaced, flag,
                                 % minimal, original
\end{lstlisting}
Every design was compiled on the full report: none produces errors or lines that overflow the
margin.

\subsection*{Before: the LaTeX default}  % the current look, for comparison
\styleinfo{original}{run-in}{67}
\setparagraphstyle{original}
\begin{sample}\sampletext\end{sample}

\subsection*{The eight designs at a glance}
\begin{center}\small\vspace{-0.6em}
\begin{tabular}{rllll}                   % number, name, key, layout, length
  \toprule
  & Design & Name & Layout & Report length \\
  \midrule
  1 & Tricolour line (your idea) & \texttt{tricolour}       & title above & 70 pages \\
  2 & Tricolour underline        & \texttt{tricolour-title} & title above & 70 pages \\
  3 & Accent line                & \texttt{accent}          & title above & 70 pages \\
  4 & Fading line                & \texttt{fade}            & title above & 70 pages \\
  5 & Label                      & \texttt{tag}             & title above & 69 pages \\
  6 & Spaced capitals            & \texttt{spaced}          & title above & 69 pages \\
  7 & Flag                       & \texttt{flag}            & run-in      & 66 pages \\
  8 & Red title                  & \texttt{minimal}         & run-in      & 66 pages \\
  \midrule
  -- & LaTeX default (now)       & \texttt{original}        & run-in      & 67 pages \\
  \bottomrule
\end{tabular}
\end{center}
{\small\emph{Title above}: the heading gets its own line. \emph{Run-in}: the text continues on
the heading's line.}

\clearpage
\showstyle{1}{Tricolour line \textnormal{\color{tpgrey}-- your idea}}{tricolour}{title above}{70}
The title in bold, with a thin line across the full text width directly under it: grey, black
and red in three equal parts, like the bar in the page footer.
\setparagraphstyle{tricolour}
\begin{sample}\sampletext\end{sample}

\showstyle{2}{Tricolour underline}{tricolour-title}{title above}{70}
The same three colours, but the line is exactly as wide as the title, so short titles get short
lines. Lighter on the page than the full-width version.
\setparagraphstyle{tricolour-title}
\begin{sample}\sampletext\end{sample}

\clearpage
\showstyle{3}{Accent line}{accent}{title above}{70}
A very light grey hairline across the page, starting with a short, thicker tricolour flag. The
colour is concentrated in one small mark, which keeps long sections calm.
\setparagraphstyle{accent}
\begin{sample}\sampletext\end{sample}

\showstyle{4}{Fading line}{fade}{title above}{70}
A red line that fades out towards the right margin. The softest of the full-width designs;
prints as a grey gradient in black and white.
\setparagraphstyle{fade}
\begin{sample}\sampletext\end{sample}

\clearpage
\showstyle{5}{Label}{tag}{title above}{69}
The title sits in a small light-grey label with a red edge, like a tab in a binder. Headings
stand out clearly when a reader scans the page.
\setparagraphstyle{tag}
\begin{sample}\sampletext\end{sample}

\showstyle{6}{Spaced capitals}{spaced}{title above}{69}
Small red capitals with extra letter spacing, followed by a thin grey line to the right margin.
An editorial look, used in magazines and annual reports.
\setparagraphstyle{spaced}
\begin{sample}\sampletext\end{sample}

\clearpage
\showstyle{7}{Flag}{flag}{run-in}{66}
A miniature of the footer flag in front of the bold title; the text continues on the same line.
Keeps the report at its current length.
\setparagraphstyle{flag}
\begin{sample}\sampletext\end{sample}

\showstyle{8}{Red title}{minimal}{run-in}{66}
The simplest change: the title in the report's red, no full stop, text on the same line. Almost
no extra space, clearly different from body text.
\setparagraphstyle{minimal}
\begin{sample}\sampletext\end{sample}

\clearpage
\section*{Notes for choosing}
\markboth{Catalogue}{}

\paragraph{Space} Designs with the title above add about one line per heading. With 101 headings
this costs 2 to 3 pages; the run-in designs keep the report at its current length (66--67
pages). The page counts above were measured on the full report.

\paragraph{Hierarchy} Sections and subsections are marked with the vertical three-colour
barettes. The paragraph designs use horizontal marks (lines, underlines, flags), so a reader
never confuses a paragraph heading with a section heading.

\paragraph{Printing} The tricolour, accent, flag and label designs stay readable in black and
white, because they rely on shape as well as colour. The fading line and the red title become
plain grey.

\paragraph{Writing titles} Write titles without a final full stop, e.g.\
\texttt{\textbackslash paragraph\{Encoder design\}}; the run-in designs add their own
separator. The 101 titles of the report have already been updated.

\paragraph{Adjusting a design} Every design is a few lines in
\texttt{setup/paragraph-styles.tex}. Common tweaks: line thickness (\texttt{0.8pt} in
\texttt{tricolour}), the gap under the title (\texttt{\textbackslash tp@lineunder}), or the
colours (\texttt{tpgrey}, \texttt{black}, \texttt{tpred}, defined in the class). This page
itself uses the \texttt{minimal} design.

\end{document}
